\documentclass[12pt]{article}
\usepackage[utf8]{inputenc}
\usepackage[T1]{fontenc}
\usepackage[margin=1in]{geometry}
\usepackage{setspace}
\usepackage{xcolor}
\usepackage[colorlinks=true, urlcolor=blue!60!black, linkcolor=blue!60!black]{hyperref}
\setlength{\parindent}{2em}
\setlength{\parskip}{1.25em}
\renewcommand{\baselinestretch}{1.0}
\begin{document}
\begin{center}
{\LARGE \textit{Regents of the University of California v. Bakke}, 438 U.S. 265 (1978)\par}
\vspace{0.8em}
{\large Case Brief --- Four-Page Edition\par}
\end{center}
\vspace{1em}

\section{Memory Jogger}

In a fractured 4--1--4 decision anchored by Justice Powell's pivotal opinion, the Court struck down UC Davis Medical School's 16-seat racial quota and ordered Allan Bakke admitted, while holding that public universities could consider race as one ``plus'' factor among many to achieve student-body diversity (\href{https://supreme.justia.com/cases/federal/us/438/265/}{Justia}).

\section{Detailed Case Facts}

UC Davis Medical School, which opened in 1968, created a two-track admissions system after its first class included no Black, Mexican-American, or American Indian students. Regular applicants with GPAs below 2.5 were rejected, and the rest were rated to produce a ``benchmark'' score. A separate special committee screened applicants who identified as disadvantaged or as members of a ``minority group''; they were exempt from the GPA cutoff, were never compared with regular applicants, and filled 16 of 100 seats. Although hundreds of disadvantaged white students applied to the special program, none was admitted.

Allan Bakke, a white engineer in his thirties with a 3.46 GPA and high MCAT scores, applied in 1973 and 1974. In 1973 he scored 468 of 500 but was rejected because no regular applicant below 470 was admitted, even though four special seats remained unfilled. In 1974 he scored 549 of 600 and was rejected again. Special admittees in both years had substantially lower averages.

\section{Procedural History}

Bakke sued under the Equal Protection Clause, the California Constitution, and Title VI. The trial court held the program an unlawful quota and barred consideration of race but did not order Bakke admitted. The California Supreme Court, applying strict scrutiny, affirmed the ban on considering race and, after the University conceded it could not prove Bakke would have been rejected anyway, ordered his admission. The U.S. Supreme Court affirmed in part and reversed in part on June 28, 1978.

\section{Judicial Votes}

A \textbf{4--1--4} split with six opinions.
\begin{itemize}
\item \textbf{Powell} announced the judgment; Brennan, White, Marshall, and Blackmun joined Parts I and V-C.
\item \textbf{Striking the quota and admitting Bakke (5--4):} Powell with Justice Stevens, Chief Justice Burger, and Justices Stewart and Rehnquist (on Title VI grounds).
\item \textbf{Lifting the ban on considering race (5--4):} Powell with Justices Brennan, White, Marshall, and Blackmun.
\item \textbf{Separate opinions:} Justices White, Marshall, and Blackmun.
\end{itemize}

\section{Holding}

Two five-Justice majorities held that (1) Title VI proscribes only racial classifications that would violate the Equal Protection Clause; (2) a rigid quota reserving seats exclusively for minority applicants is unlawful, and Bakke must be admitted because the University could not prove he would have been rejected anyway; and (3) a university is not barred from considering race as one factor. No other Justice joined Powell's view that student-body diversity is a compelling interest; \textit{Grutter v. Bollinger} (2003) later endorsed it, before \textit{SFFA v. Harvard} (2023) effectively repudiated race-conscious admissions.

\section{Analysis of Opinions}

\textbf{Powell.} Title VI incorporates the Equal Protection standard. Equal protection rights are ``personal rights,'' and in a ``Nation of minorities'' courts have no baseline for ranking which groups deserve preference, so \textit{all} racial classifications get strict scrutiny. Powell rejected three of Davis's interests: reducing the deficit of minorities in medicine (``discrimination for its own sake''), countering ``societal discrimination'' absent findings of specific violations, and improving care in underserved areas (no proof). He accepted a fourth, a diverse student body, grounded in academic freedom and the ``robust exchange of ideas.'' The two-track quota was unnecessary because it insulated minority candidates from competition; the Harvard plan, treating race as a ``plus'' in individualized review, was permissible.

\textbf{Brennan, White, Marshall, and Blackmun.} Colorblindness is an aspiration, not reality. Benign remedial classifications should get \textbf{intermediate scrutiny}, which Davis's program satisfied as a remedy for past societal discrimination; they saw no constitutional difference between Davis's set-aside and Harvard's ``plus.''

\textbf{White, Marshall, and Blackmun (separately).} White found no private right of action under Title VI. Marshall traced slavery, \textit{Plessy}, and Jim Crow, noting the irony of using the Fourteenth Amendment to bar class-based remedies. Blackmun noted that universities have always favored athletes, alumni children, and others.

\textbf{Stevens, Burger, Stewart, and Rehnquist.} Avoiding the Constitution, they read Title VI as ``crystal clear'': a federally funded program may not exclude anyone because of race.

\section{Examples: Future Applications}

\begin{enumerate}
\item \textbf{Same-Side (Invalid): Automatic 20-Point Bonus.} An index awarding 20 points to every underrepresented minority applicant, as in \textit{Gratz v. Bollinger}, makes race mechanical rather than individualized.
\item \textbf{Same-Side (Invalid): Minority-Only Track.} A law school reserves 15\% of seats for a subcommittee open only to Black and Hispanic applicants, replicating Davis's two tracks.
\item \textbf{Opposite-Side (Valid): Holistic ``Plus'' Review.} A single pool weighing race as a flexible factor with no numerical target tracks the Harvard plan (and \textit{Grutter}, later curtailed by \textit{SFFA}).
\item \textbf{Opposite-Side (Valid): Socioeconomic Preference.} A ``plus'' for low-income and first-generation applicants of any race involves no racial classification.
\item \textbf{Fence-Sitter: Top 10\% Plus Holistic Review.} Race as a modest tiebreaker after a race-neutral percentage plan, as in \textit{Fisher} (2013), may fill gaps or may prove unnecessary.
\end{enumerate}

\section{Critique}

\textbf{Originalist critique.} In \textit{The Disease as Cure} (1979), then-Professor Antonin Scalia argued that racial preferences perpetuate resentment, stigmatize minority achievers, and violate equal individual treatment. Textualists echoed Stevens's reading of Title VI and called Powell's quota-versus-plus line a semantic veil, a critique that prevailed in \textit{SFFA}.

\textbf{Progressive critique.} Antisubordination scholars call Powell's opinion a loss for racial justice: by rejecting ``societal discrimination'' and applying strict scrutiny to remedial measures, it shifted Equal Protection toward symmetrical anticlassification, treated minority students as educational assets for white peers, and left legacy preferences intact.

\section{Key Quotations}

\begin{quote}
``The guarantee of equal protection cannot mean one thing when applied to one individual and something else when applied to a person of another color. If both are not accorded the same protection, then it is not equal.''
\end{quote}
\noindent --- Justice Powell, announcing the judgment of the Court, 438 U.S. at 289--290.

\begin{quote}
``Preferring members of any one group for no reason other than race or ethnic origin is discrimination for its own sake. This the Constitution forbids.''
\end{quote}
\noindent --- Justice Powell, 438 U.S. at 307.

\begin{quote}
``In such an admissions program, race or ethnic background may be deemed a `plus' in a particular applicant's file, yet it does not insulate the individual from comparison with all other candidates for the available seats.''
\end{quote}
\noindent --- Justice Powell, 438 U.S. at 317.

\end{document}
